\section{The Harness: Lean or Heavy?}\label{sec:harness}

The harness is everything configured around the model: tools, context files, skills, sub-agents, hooks, permissions and memory. A source-level reading of Claude Code finds ``a simple while-loop'' surrounded by a permission system with seven modes, a five-layer compaction pipeline and four extensibility mechanisms \citep{H16}. The evidence, mostly from short benchmark tasks rather than long unattended runs, favors a lean harness in which each component earns its place, and it places the controls that matter for safety in deterministic code outside the model.

\subsection{What Helps and What Often Does Not}

The components with measured benefits are those that shape what the agent sees and does at each step. In the SWE-agent study, a purpose-built tool interface solved 10.7 percentage points more SWE-bench Lite instances than a plain shell \citep{H10}. In SkillsBench, curated skills raised the average pass rate from 33.9\% to 50.5\%, and focused skills with at most three modules beat larger bundles \citep{H6}. Simple context management matches clever context management: in a JetBrains Research study, dropping old tool outputs ``halves cost'' relative to an unmanaged agent while matching LLM summarization on SWE-bench Verified \citep{H3}, and learned compression cut peak tokens by 26--54\% \citep{H4}. Rewriting a context wholesale causes ``context collapse,'' which incremental curation avoids \citep{H5}. Input length alone degrades performance by 13.9\% to 85\% even with perfect retrieval \citep{F22}.

The heavier, more popular components fare worse. An ETH Zurich study found that repository context files (\texttt{AGENTS.md}, \texttt{CLAUDE.md}) do not raise success rates on average while increasing inference cost by over 20\% when the files are LLM-generated; developer-written files raised success by a non-significant 2.4\% on average, helping every agent tested except Claude Code, at up to 19\% more cost \citep{H1}. Yet context files appear in 2,586 of 2,853 repositories surveyed, against 158 with skills and 131 with sub-agents \citep{H2}. Skills the agent writes for itself land below the no-skill baseline, by 8.1 points for Claude Code with Opus 4.7 \citep{H6}, and gains from curated skills shrink toward that baseline when the agent must find the right skill in a large realistic pool \citep{H7}. Skills also carry risk and upkeep: 26.1\% of 31,132 marketplace skills contained at least one vulnerability, with script-bundling skills 2.12 times as likely to \citep{H22}, and their instructions and scripts must be maintained together \citep{H12}. Harness gains can decay as models improve: in a Meta and Harvard study, context management added 6.6 points for Claude 4 Sonnet on a SWE-Bench-Pro subset but less for Claude 4.5 \citep{H20}. In practice, a user who installed every agent in one large marketplace received a warning of about 404.1k tokens of agent descriptions \repo{wshobson/agents}{93}, and a workflow skill re-attached after compaction was cut at 20,000 characters, dropping its five-round cap; one task then ran 18 fix rounds, although the reporter notes that this evidence does not establish the cause \repo{obra/superpowers}{2481}.

\subsection{Permissions and Containment}

Unattended runs remove the human who approves risky actions, and the evidence favors deterministic controls over model judgment. Anthropic reports that users accept 93\% of permission prompts, and that its auto-mode classifier has a 0.4\% false-positive rate on internal traffic but a 17\% false-negative rate on 52 real overeager actions \citep{H23}. An independent stress test of the same mode found an 81.0\% end-to-end false-negative rate (95\% CI 73.8--87.4\%) on 128 deliberately ambiguous DevOps prompts; 36.8\% of state-changing actions were in-project file edits that never reached the classifier, and on the actions it did evaluate the rate was still 70.3\%; the authors caution against extrapolating to production traffic \citep{H17}. Deterministic policies do better on their benchmarks: rules checked before each action prevented unsafe executions in over 90\% of code-agent cases with millisecond overhead \citep{H8}, least-privilege policies cut attack success on AgentDojo from 39.9\% to 1.0\% \citep{H9}, and CaMeL solved 77\% of AgentDojo tasks with provable security against 84\% undefended \citep{H18}. Anthropic reports that operating-system sandboxing cut permission prompts by 84\% in internal use \citep{H24}. Containment only holds if the agent cannot lift it: an open issue reports that, in auto-allow mode, Claude can disable the sandbox by setting \texttt{dangerouslyDisableSandbox} after a sandbox-caused failure \repo{anthropics/sandbox-runtime}{97}, and a since-closed issue reported that one plugin's safety hooks ``appear installed and simply never fire'' on Windows \repo{uditgoenka/autoresearch}{115}. The stronger designs keep the agent away from its own controls, for example by restoring Claude configuration files from the base branch before a run on a pull request \repofile{anthropics/claude-code-action}{docs/security.md}.

\subsection{State Across Context Resets}

Long runs outlive any context window, so progress must live outside it. Anthropic's recipe for long-running agents keeps a progress file alongside the git history, and reports that working on one feature at a time ``turned out to be critical'' \citep{H13}. Measured gains are modest: rerunning 151 tasks with notes from a first run raised the resolve rate from 53\% to 54.4\% \citep{H20}. File-based state has its own failure modes. A repository-global plan pointer let one session load another session's plan after compaction \repo{OthmanAdi/planning-with-files}{240}; the same project warns that ``a checked box is a progress claim, not evidence'' \repofile{OthmanAdi/planning-with-files}{docs/attestation-locking.md}; and a design proposal in another framework notes that a sub-agent ``that ran off the rails but returns `completed' has its conclusions adopted wholesale'' \repo{bytedance/deer-flow}{4651}. Durable state solves memory, not verification: a plan file records what the agent claims, and something else must check it. For builders of such systems, the evidence suggests adding harness only where a paired run at a fixed model shows it helps.
